\subsection{Sextic tensor decomposition for formaldehyde and thioformaldehyde}

For the sextic part of the tensor analysis it is convenient to consider
small asymmetric tops for which complete sets of sextic centrifugal
distortion constants are available both from experiment and from direct
rovibrational calculations.  Formaldehyde (H$_2$CO) and thioformaldehyde
(H$_2$CS) provide particularly suitable examples in this respect, because
both molecules were analyzed in Ref.~\citenum{Dinu2022} using Watson's
$S$-reduced Hamiltonian in the $I^r$ representation, and both experimental
and variationally fitted sextic constants were reported.

These two molecules therefore allow one to illustrate directly the
decomposition
\[
\mathbf H_S = \mathbf B_{I^r}\,\boldsymbol\sigma + \mathbf r,
\]
introduced above.  Here $\boldsymbol\sigma$ is the five-dimensional vector
of physical sextic coordinates on the canonical invariant subspace, whereas
$\mathbf r$ measures the residual component in the complementary
two-dimensional subspace.  In the present implementation the sextic
representation transforms act on $\boldsymbol\sigma$, while $\mathbf r$
provides a diagnostic of the consistency of the input parameter set with the
modeled physical subspace.

The resulting coordinates for H$_2$CO and H$_2$CS are collected in
Table~\ref{tab:sextic_sigma_h2co_h2cs}.  For each molecule two sets of
sextic constants were considered: the experimentally derived constants and
the corresponding VCI/PFIT constants reported in Ref.~\citenum{Dinu2022}.
The values of the five canonical coordinates $\sigma_1,\ldots,\sigma_5$ are
reported together with the residual norms
\[
r_{\max}=\max_i |r_i|,
\qquad
\|r\|_2 = \left(\sum_i r_i^2\right)^{1/2}.
\]

\begin{table*}[ht]
\centering
\caption{Canonical sextic coordinates $\boldsymbol\sigma$ and residual norms
for H$_2$CO and H$_2$CS obtained from the sextic Watson constants reported in
Ref.~\citenum{Dinu2022}.  All quantities are given in kHz after conversion of
the tabulated sextic constants to a common MHz unit and subsequent tensor
decomposition in the $I^r$ representation and $S$ reduction.}

\begin{tabular}{llrrrrrrr}
\hline
Molecule & Source & $\sigma_1$ & $\sigma_2$ & $\sigma_3$ & $\sigma_4$ & $\sigma_5$ & $r_{\max}$ & $\|r\|_2$ \\
\hline
H$_2$CO & Experiment & -0.8542 & -0.0470 & 3.2566 & -0.1677 & 2.4863 & 1.0219 & 1.7596 \\
H$_2$CO & VCI/PFIT   & -0.8804 & -0.0468 & 3.3571 & -0.1729 & 2.5632 & 1.0534 & 1.8135 \\
H$_2$CS & Experiment & -1.2630 & -0.0826 & 4.8065 & -0.2471 & 3.6670 & 1.5101 & 2.6028 \\
H$_2$CS & VCI/PFIT   & -1.2775 & -0.0824 & 4.8615 & -0.2500 & 3.7090 & 1.5275 & 2.6325 \\
\hline
\end{tabular}

\label{tab:sextic_sigma_h2co_h2cs}
\end{table*}

Several useful observations follow from these data.  First, for each
molecule the experimental and VCI/PFIT coordinates are very similar.  This
shows that the canonical vector $\boldsymbol\sigma$ provides a compact and
stable representation of the physical sextic content, even though the
individual Watson constants may differ noticeably between experiment and
theory.  Second, H$_2$CO and H$_2$CS display the same qualitative pattern in
their sextic coordinates, with H$_2$CS showing a systematically larger
magnitude for all five components.  This is consistent with the larger
overall sextic centrifugal distortion expected for the heavier sulfur
analogue.

The residual norms are not zero for the tabulated parameter sets.
This is not a failure of the tensor algorithm.  Rather, it indicates that
the published sextic constants, when taken as finite-precision Watson
parameter sets, do not lie exactly on the modeled five-dimensional subspace.
This behaviour is expected for real fitted parameter sets, which are affected
by parameter correlations, finite numerical precision, and differences in
fitting strategy.  The residual therefore plays the role of a diagnostic
quantity: it measures the distance of a given sextic parameter set from the
canonical tensor model adopted in the present work.

Once the sextic constants are projected onto the physical subspace, the
resulting vector $\boldsymbol\sigma$ is transported unchanged by the
representation transformation algorithm.  In this sense, the sextic
decomposition provides the direct analogue of the quartic tensor analysis:
the Watson constants are representation-dependent coordinates, whereas the
canonical physical coordinates $\boldsymbol\sigma$ provide the
representation-independent sextic content.  In practical terms, this means
that sextic parameter sets obtained in different representations may be
converted to the same canonical coordinates, compared together with their
residual norms and transformation condition numbers, and only then used as
input for any final refit in the representation judged to be numerically most
favourable.
